\begin{exercises}
  \item[]\textit{这里总结本书中 $2$-\hbox{} 与 $3$-置换所用的记号。
    \begin{center}
      \begin{tabular}[t]{c|cc}
        $i$          &$1$      &$2$    \\
        \hline
        $\phi_1(i)$  &$1$      &$2$     \\
        $\phi_2(i)$  &$2$      &$1$     
      \end{tabular}
      \qquad
      \begin{tabular}[t]{c|ccc}
        $i$          &$1$     &$2$   &$3$    \\
        \hline
        $\phi_1(i)$  &$1$     &$2$   &$3$    \\
        $\phi_2(i)$  &$1$     &$3$   &$2$    \\
        $\phi_3(i)$  &$2$     &$1$   &$3$    \\
        $\phi_4(i)$  &$2$     &$3$   &$1$    \\
        $\phi_5(i)$  &$3$     &$1$   &$2$    \\
        $\phi_6(i)$  &$3$     &$2$   &$1$    
      \end{tabular}
    \end{center}  }
  \item 
    写出一般 $\nbyn{2}$ 矩阵及其转置的置换展开。
    \begin{answer}
      这是 $\nbyn{2}$ 矩阵的行列式的置换展开
      \begin{equation*}
        \begin{vmat}
          a  &b  \\
          c  &d
        \end{vmat}
         =ad\cdot \begin{vmat}[r]
                    1  &0  \\
                    0  &1
                  \end{vmat}
         +bc\cdot \begin{vmat}[r]
                    0  &1 \\
                    1  &0
                  \end{vmat}
      \end{equation*}
      以及它的转置的行列式的置换展开。
      \begin{equation*}
         \begin{vmat}
            a  &c  \\
            b  &d
         \end{vmat}
         =ad\cdot\begin{vmat}[r]
                    1  &0  \\
                    0  &1
                 \end{vmat}
         +cb\cdot\begin{vmat}[r]
                   0  &1 \\
                   1  &0
                 \end{vmat}
       \end{equation*}
       与小节中描述的 $\nbyn{3}$ 展开一样，对应项的置换矩阵互为转置
       （尽管每个矩阵都是自转置的这一事实把这一点掩盖了）。
    \end{answer}
  \recommended \item 
    \textit{本题也出现在前面的小节中。}
    \begin{exparts}
      \partsitem 求每个 $2$-置换的逆。
      \partsitem 求每个 $3$-置换的逆。
    \end{exparts}
    \begin{answer}
      每一个都容易验证。
      \begin{exparts}
        \partsitem 
          \begin{tabular}[t]{r|cc}
            \textit{置换} &$\phi_1$  &$\phi_2$ \\
             \hline
            \textit{逆}     &$\phi_1$  &$\phi_2$ 
          \end{tabular}
        \partsitem 
          \begin{tabular}[t]{r|cccccc}
            \textit{置换} 
              &$\phi_1$ &$\phi_2$ &$\phi_3$ &$\phi_4$ &$\phi_5$ &$\phi_6$ \\
            \hline
            \textit{逆}
              &$\phi_1$ &$\phi_2$ &$\phi_3$ &$\phi_5$ &$\phi_4$ &$\phi_6$ 
          \end{tabular}
      \end{exparts}
    \end{answer}
  \recommended \item
    \begin{exparts}
      \partsitem 求每个 $2$-置换的符号。
      \partsitem 求每个 $3$-置换的符号。
    \end{exparts}
    \begin{answer}
      \begin{exparts}
        \partsitem \( \sgn(\phi_1)=+1 \), \( \sgn(\phi_2)=-1 \)
        \partsitem \( \sgn(\phi_1)=+1 \), \( \sgn(\phi_2)=-1 \),
              \( \sgn(\phi_3)=-1 \), \( \sgn(\phi_4)=+1 \),
              \( \sgn(\phi_5)=+1 \), \( \sgn(\phi_6)=-1 \)
      \end{exparts}  
     \end{answer}
  \item 
    找出此矩阵的置换展开中唯一非零的项。
    \begin{equation*}
      \begin{vmat}[r]
        0  &1  &0  &0  \\
        1  &0  &1  &0  \\
        0  &1  &0  &1  \\
        0  &0  &1  &0
      \end{vmat}
    \end{equation*}
    通过求相应置换的符号来计算该行列式。
    \begin{answer}
      要在置换展开中得到非零项，我们必须取 \( 1,2 \) 处的元素与
      \( 4,3 \) 处的元素。
      选定这两处之后，我们还必须取 \( 2,1 \) 处的元素与 \( 3,4 \) 处的元素。
      \( \sequence{2,1,4,3} \) 的符号是 \( +1 \)，因为由
      \begin{equation*}
        \begin{mat}[r]
          0  &1  &0  &0  \\
          1  &0  &0  &0  \\
          0  &0  &0  &1  \\
          0  &0  &1  &0            
        \end{mat}
      \end{equation*}
      经过两次行对换 $\rho_1\leftrightarrow\rho_2$ 与
      $\rho_3\leftrightarrow\rho_4$ 即可得到单位矩阵。  
    \end{answer}
  \item  
    \cite{Strang}
    $n$-置换 \( \phi=\sequence{n,n-1,\dots,2,1} \) 的符号是什么？
    \begin{answer}
      规律如下。
      \begin{center}
         \begin{tabular}[b]{c|ccccccc}
           \( i \) &\( 1 \) &\( 2 \) &\( 3 \) &\( 4 \) &\( 5 \)
              &\( 6 \) &\ldots \\
           \hline
           \( \sgn(\phi_i) \) &\( +1 \) &\( -1 \) &\( -1 \) &\( +1 \)
              &\( +1 \) &\( -1 \) &\ldots
         \end{tabular}
      \end{center}  
      于是为求 $\phi_{n!}$ 的符号，我们用 $n!-1$ 减一，再看它除以四的余数。
      若余数是 $1$ 或 $2$，则符号为 $-1$；否则为 $+1$。
      对 $n>4$，数 $n!$ 能被四整除，所以 $n!-1$ 除以四的余数是
      $-1$（更准确地说，余数是 $3$），因此符号是 $+1$。
      $n=1$ 的情形符号为 $+1$，$n=2$ 的情形符号为 $-1$，
      $n=3$ 的情形符号为 $-1$。
    \end{answer}
  \item \label{exer:PermInvFacts}
     证明这些结论。   
     \begin{exparts}
        \partsitem 每个置换都有逆。
        \partsitem \( \sgn(\phi^{-1})=\sgn(\phi) \)
        \partsitem 每个置换都是另一个置换的逆。
     \end{exparts}
    \begin{answer}
      \begin{exparts}
       \partsitem 我们可以把置换看成
         \( \map{\phi}{\set{1,\ldots,n}}{\set{1,\ldots,n}} \)
         这样单射且满射的映射。
         任何单射且满射的映射都有逆。
       \partsitem 若从 \( P_\phi \)
         变到单位矩阵总需要奇数次对换，那么从单位矩阵变到
         \( P_\phi \) 也总需要奇数次对换（任何对换都是可逆的）。
       \partsitem 这就是第一问。
      \end{exparts}  
    \end{answer}
  \item \label{exer:PermInvIffMatTrans}
      证明对任意置换 $\phi$，置换逆的矩阵就是置换的矩阵的转置
      $P_{\phi^{-1}}=\trans{P_{\phi}}$。
      \begin{answer}
        若 $\phi(i)=j$ 则 $\phi^{-1}(j)=i$。
        于是注意到 $P_{\phi}$ 在 $i,j$ 处有一个 $1$ 当且仅当 $\phi(i)=j$，
        而 $P_{\phi^{-1}}$ 在 $j,i$ 处有一个 $1$ 当且仅当 $\phi^{-1}(j)=i$，
        结论即得。
      \end{answer}
  \recommended \item 
    证明有 \( m \) 个逆序的置换矩阵可以经过 \( m \) 步行对换变成单位矩阵。
    将此与\nearbycorollary{cor:ParityInversEqParitySwaps}对比。
    \begin{answer}
      这并不是说 \( m \) 是变成单位矩阵所需的最少对换次数，
      也不是说 \( m \) 是最多的。
      它说的是存在一种恰好经 \( m \) 步对换变成单位矩阵的方法。

      设 \( \iota_j \) 是相对于前面某行来说第一个出现逆序的行，
      并设 \( \iota_k \) 是给出该逆序的第一行。
      我们有这样一个行区间。
      \begin{equation*}
         \begin{mat}
           \vdots      \\
           \iota_k     \\
           \iota_{r_1} \\
           \vdots      \\
           \iota_{r_s} \\
           \iota_j     \\
           \vdots
         \end{mat}
         \qquad j<k<r_1<\dots <r_s
      \end{equation*}
      作对换。
      \begin{equation*}
         \begin{mat}
           \vdots      \\
           \iota_j     \\
           \iota_{r_1} \\
           \vdots      \\
           \iota_{r_s} \\
           \iota_k     \\
           \vdots
         \end{mat}
      \end{equation*}
      第二个矩阵的逆序少了一个，因为该区间内的逆序少了一个
      （\( s \) 对 \( s+1 \)），而涉及区间之外各行的逆序不受影响。

      如此进行下去，每作一次行对换就把逆序的个数减少一个。
      当逆序不复存在时，结果就是单位矩阵。

      与\nearbycorollary{cor:ParityInversEqParitySwaps}的对比在于，
      本练习的陈述是一个“存在”命题：
      存在一种恰好经~$m$ 步对换变成单位矩阵的方法。
      而该推论是一个“对所有”命题：对变成单位矩阵的所有对换方式，
      奇偶性（是偶还是奇）都相同。 
    \end{answer}
  \recommended \item 
    对任意置换 $\phi$，令 $g(\phi)$ 是如下定义的整数。
    \begin{equation*}
       g(\phi)=\prod_{i<j} [\phi(j)-\phi(i)]
    \end{equation*}
    （这是在所有满足 \( i<j \) 的指标 \( i \) 与 \( j \) 上，
    对上述形式的项作的乘积。）
    \begin{exparts}
      \partsitem 计算 $g$ 在所有 $2$-置换上的值。
      \partsitem 计算 $g$ 在所有 $3$-置换上的值。
      \partsitem 证明 $g(\phi)$ 不为 $0$。
      \partsitem 证明这个结论。
         \begin{equation*}
            \sgn(\phi)=\frac{g(\phi)}{\absval{g(\phi)}}
         \end{equation*}
    \end{exparts}
    许多作者把这个公式当作符号函数的定义。
    \begin{answer}
      \begin{exparts}
        \partsitem 首先，$g(\phi_1)$ 是单独一个因子 $2-1$ 的乘积，
          所以 $g(\phi_1)=1$。
          其次，$g(\phi_2)$ 是单独一个因子 $1-2$ 的乘积，
          所以 $g(\phi_2)=-1$。
        \partsitem 
            \begin{tabular}[t]{r|cccccc}
               \textit{置换 $\phi$} 
                 &$\phi_{1}$ &$\phi_{2}$ &$\phi_{3}$ 
                    &$\phi_{4}$ &$\phi_{5}$ &$\phi_{6}$ \\
               \hline
               $g(\phi)$ &$2$ &$-2$ &$-2$ &$2$ &$2$ &$-2$ 
            \end{tabular}
        \partsitem 它是一些非零项的乘积。
        \partsitem 注意 \( \phi(j)-\phi(i) \) 为负当且仅当
           \( \iota_{\phi(j)} \) 与 \( \iota_{\phi(i)} \) 处于它们通常顺序的一个逆序之中。
      \end{exparts}  
     \end{answer}
%  \item Prove that the signum function on permutations is unique in that
%    if \( f \) is a map from the permutations of \( \set{1,\ldots,n} \)
%    to the integers that is multiplicative:
%    \( f(\sigma\phi)=f(\sigma)f(\phi) \) then \( f \) is either the
%    function that always gives zero, or the function that always gives
%    one, or is the signum.
%    \begin{answer}
%      Consider the identity permutation.
%      \begin{equation*}
%        \map{\identity}{\set{1,2,\ldots,n}}{\set{1,2,\ldots,n}}    
%      \end{equation*}
%      Note that $f(\identity)$ equals $0$ or $1$ because
%      the multiplicative property gives 
%      $f(\identity)
%       =f(\composed{\identity}{\identity})
%       =f(\identity)\cdot f(\identity)
%       =f(\identity)^2$.
%      If $f(\identity)=0$ then $f$ is the constant function 
%      $f(\phi)=0$ because 
%      $f(\phi)=f(\composed{\identity}{\phi})=0\cdot f(\phi)=0$.
%
%      If $f(\identity)=1$ then consider the permutation
%      $\tau$ that just swaps $i$ and $j$
%      \begin{equation*}
%        1\mapsto 1
%        \quad 2\mapsto 2
%        \quad \cdots 
%        \quad i\mapsto j
%        \quad \cdots
%        \quad j\mapsto i
%        \quad \cdots
%        \quad n\mapsto n
%      \end{equation*}
%      (the $n=1$ special case is easy).
%      Note that $\composed{\tau}{\tau}=\identity$ so 
%      $1=f(\composed{\tau}{\tau})=f(\tau)\cdot f(\tau)$,
%      and $f(\tau)$ is either $+1$ or $-1$.
%    \end{answer}
\end{exercises}



\index{determinant|)}
